\section{La evidencia de comportamiento}\label{sec:evidencia}

El contador de bits es exacto para lo que sabe leer y ciego para lo demás. La evidencia de
comportamiento es la segunda línea: no mira \emph{qué} sale sino \emph{cómo se ve} la llamada. Es un
modelo estadístico pequeño, ajustado con datos, y es la pieza que frenó las familias nuevas del set de
prueba. Está en \cod{engine/evidence/} y en \cod{engine/detector/features.py}.

Esta sección lo describe entero, incluido lo que el ajuste dejó a la vista: el modelo quedó saturado y
se comporta casi como una regla de dos valores. Decirlo es más útil que presentarlo como algo más fino
de lo que es.

\subsection{La idea}

A cada llamada se le asigna un número $\ell_t$, positivo si la llamada se parece más a las de un
ataque y negativo si se parece más al trabajo normal. Esos números se acumulan en una suma $S_t$ que
no baja de cero. Cuando $S_t$ pasa un umbral, hay evidencia de que la sesión cambió de régimen. Es el
esquema de detección de cambios de Page~\cite{page}, conocido como CUSUM, aplicado a una razón de
verosimilitudes.

\subsection{Las nueve señales}

De cada llamada se extraen nueve números (\cref{tab:senales}). Tres vienen del grafo de archivos,
tres del grafo de llamadas, uno del proxy y dos del contenido.

\begin{table}[ht]
\caption{Las nueve señales de una llamada (\cod{engine/detector/features.py}).}\label{tab:senales}
\small
\begin{tabularx}{\textwidth}{@{}P{3.7cm}P{1.9cm}L@{}}
\toprule
Señal & Origen & Qué mide \\
\midrule
$m$ & Grafo de archivos & Sensibilidad del nodo más sensible que toca la llamada \\
$m_{\mathrm{sd}}$ & Grafo de archivos & Incertidumbre de esa sensibilidad \\
\cod{jump} & Grafo de archivos & Saltos entre lo que toca esta llamada y lo que tocaron las cinco
  anteriores \\
$U$ & Grafo de llamadas & 1 si la llamada es no confiable o desciende de una que lo es \\
$n_S$ & Grafo de llamadas & Número de claves que la llamada leyó o heredó \\
\cod{is_sink} & Proxy & 1 si la llamada saca datos \\
\cod{prev_source_now_sink} & Grafo de llamadas & 1 si la llamada anterior fue una fuente y esta es un
  sumidero \\
\cod{arg_entropy} & Contenido & Entropía de Shannon, por carácter, de los argumentos serializados \\
\cod{short_match} & Contador & Bits candidatos vigentes, sumados sobre todas las claves \\
\bottomrule
\end{tabularx}
\end{table}

\subsection{El modelo}

El modelo (\cod{engine/evidence/model.py}) trata cada señal por separado. Para la señal $j$ divide su
rango en cajas, con bordes en los cuantiles de las muestras normales y de ataque juntas (hasta 10
cajas; los bordes repetidos se funden). En cada caja $c$ estima dos probabilidades, la de las llamadas
normales y la de las de ataque, con suavizado de Laplace $\alpha = 1$:
\begin{equation}
  \hat p_0^{\,j}(c) = \frac{N_0^{\,j}(c) + \alpha}{N_0 + \alpha\, k_j} , \qquad
  \hat p_1^{\,j}(c) = \frac{N_1^{\,j}(c) + \alpha}{N_1 + \alpha\, k_j} ,
  \label{eq:hist}
\end{equation}
donde $N_0^{\,j}(c)$ y $N_1^{\,j}(c)$ cuentan las muestras normales y de ataque en la caja, $N_0$ y
$N_1$ son los totales y $k_j$ es el número de cajas. El incremento de una llamada con señales
$f_{t1}, \dots, f_{t9}$ es la suma de los logaritmos de las razones, recortada a $[-B, B]$ con
$B = 3$:
\begin{equation}
  \ell_t = \clip\!\left( \sum_{j=1}^{9} \log \frac{\hat p_1^{\,j}\!\left(c_j(f_{tj})\right)}
                                               {\hat p_0^{\,j}\!\left(c_j(f_{tj})\right)} ,\; -B ,\; B \right) .
  \label{eq:ell}
\end{equation}
Sumar los logaritmos equivale a suponer que las señales son independientes dada la clase. No lo son:
varias dicen casi lo mismo. El recorte existe para que ninguna llamada, por muy segura que parezca,
mueva la suma más de $B$.

\subsection{La acumulación}

La evidencia acumulada sigue la recurrencia de Page,
\begin{equation}
  S_t = \max\!\left(0,\; S_{t-1} + \ell_t\right) , \qquad S_0 = 0 .
  \label{eq:cusum}
\end{equation}
Hay una suma por sesión. No baja de cero, de modo que una sesión larga de trabajo normal no acumula
«crédito» que después tape un ataque. No se reinicia cuando una clave rota.

Dos umbrales la leen. Con $S_t \geq \brot$, y si hay una clave expuesta, se rota. Con $S_t > b$, y si
además hay un camino de contenido en el grafo de llamadas, se bloquea (\cref{sec:decision}).

\subsection{El ajuste}\label{sec:modelo-ajustado}

El modelo se ajusta con \cod{python -m eval.fit}. Las señales se obtienen pasando cada sesión por el
detector en modo de solo medición y sin modelo.

\begin{itemize}
  \item \textbf{Llamadas normales:} las \num{2511} llamadas de las 60 sesiones de entrenamiento del
    tráfico normal de fondo (la mitad de las 120 de \cod{data/benign}).
  \item \textbf{Llamadas de ataque:} los 458 envíos etiquetados de los 80 ataques de
    \cod{redteam/train}. El resto de las llamadas de una sesión de ataque no entra, ni como normal ni
    como ataque.
\end{itemize}

Ni validación ni prueba entran al ajuste. La \cref{fig:modelo} muestra el modelo completo: las nueve
señales, caja por caja. El \apref{ap:modelo} da los bordes y los valores.

\begin{figure}[ht]
\centering
\begin{tikzpicture}
\begin{groupplot}[
  group style={group size=3 by 3,horizontal sep=1.25cm,vertical sep=1.25cm},
  width=5.5cm,height=3.55cm,ymin=-7.2,ymax=9,ytick={-6,-3,0,3,6},
  xticklabel style={font=\tiny},tick label style={font=\tiny},
  extra y ticks={-3,3},extra y tick labels={},
  extra y tick style={grid=major,grid style={oro,dashed,line width=0.5pt}},
  axis x line=bottom,axis y line=left,enlarge x limits=false]
\nextgroupplot[title={$U$}, xmin=0.3, xmax=2.7, xtick={1,...,2}]
\addplot[ybar, bar width=8pt, draw=none, fill=tinta] coordinates {(1,-6.133) (2,7.825)};
\nextgroupplot[title={\texttt{is\_sink}}, xmin=0.3, xmax=2.7, xtick={1,...,2}]
\addplot[ybar, bar width=8pt, draw=none, fill=tinta] coordinates {(1,-6.123) (2,4.567)};
\nextgroupplot[title={\texttt{short\_match}}, xmin=0.3, xmax=2.7, xtick={1,...,2}]
\addplot[ybar, bar width=8pt, draw=none, fill=tinta] coordinates {(1,-1.494) (2,5.271)};
\nextgroupplot[title={\texttt{arg\_entropy}}, xmin=0.3, xmax=10.7, xtick={1,...,10}]
\addplot[ybar, bar width=5pt, draw=none, fill=tinta] coordinates {(1,-3.954) (2,-3.972) (3,-4.091) (4,-3.989) (5,-4.030) (6,-3.976) (7,-4.049) (8,-0.456) (9,4.435) (10,1.650)};
\nextgroupplot[title={$m$}, xmin=0.3, xmax=7.7, xtick={1,...,7}]
\addplot[ybar, bar width=5pt, draw=none, fill=tinta] coordinates {(1,1.360) (2,-3.996) (3,-4.013) (4,-4.165) (5,-4.096) (6,-2.487) (7,-4.622)};
\nextgroupplot[title={$m_{\mathrm{sd}}$}, xmin=0.3, xmax=7.7, xtick={1,...,7}]
\addplot[ybar, bar width=5pt, draw=none, fill=tinta] coordinates {(1,1.284) (2,-4.026) (3,-3.793) (4,-4.026) (5,-3.972) (6,-4.096) (7,-4.234)};
\nextgroupplot[title={\texttt{jump}}, xmin=0.3, xmax=4.7, xtick={1,...,4}]
\addplot[ybar, bar width=8pt, draw=none, fill=tinta] coordinates {(1,0.647) (2,-4.408) (3,-4.385) (4,-4.097)};
\nextgroupplot[title={$n_S$}, xmin=0.3, xmax=1.7, xtick={1,...,1}]
\addplot[ybar, bar width=8pt, draw=none, fill=tinta] coordinates {(1,-0.002)};
\nextgroupplot[title={\texttt{prev\_source\_now\_sink}}, xmin=0.3, xmax=1.7, xtick={1,...,1}]
\addplot[ybar, bar width=8pt, draw=none, fill=tinta] coordinates {(1,-0.002)};

\end{groupplot}
\end{tikzpicture}
\caption{El modelo de evidencia ajustado (\cod{models/model.json}): $\log(\hat p_1/\hat p_0)$ de cada
caja de cada señal, con las cajas en orden creciente del valor de la señal. Las líneas punteadas
marcan el tope $\pm B = \pm 3$ del incremento. Varias señales, solas, ya lo superan. Las dos
últimas tienen una sola caja con datos, con un valor casi nulo.}
\label{fig:modelo}
\end{figure}

El informe del ajuste (\cod{models/model.report.json}) dice qué tan separadas quedaron las dos clases:

\begin{itemize}
  \item el \qty{99.92}{\percent} de las llamadas normales de entrenamiento queda exactamente en el
    tope inferior, $\ell = -3$;
  \item la media de $\ell$ en las llamadas normales es \num{-2.9995} y su varianza, \num{0.0003};
  \item la media de $\ell$ en los envíos de ataque es \num{3.0}: todos quedan en el tope superior.
\end{itemize}

\begin{clave}
\textbf{El modelo está saturado.} En los datos de entrenamiento no produce valores intermedios: da
$-3$ o $+3$. Funciona como un clasificador de dos salidas, no como una medida graduada de sospecha. La
razón está en la \cref{fig:modelo}: la suma junta varias señales que valen cada una más que el tope, y
que además dicen casi lo mismo (si la llamada saca datos, si desciende de algo no confiable).
\end{clave}

Dos de las nueve señales, $n_S$ y \cod{prev_source_now_sink}, quedaron sin información: en los datos
de ajuste casi todas sus muestras valen 0, los bordes se funden en uno solo y todas caen en la misma
caja, con una razón cercana a 1.

\subsection{Qué aprendió, en palabras}

Para saber qué regla quedó escrita en esos histogramas, medimos $\ell$ en cada una de las
\num{15281} llamadas del set de prueba, con las marcas del proxy (\cref{tab:ell}).

\begin{table}[ht]
\caption{El incremento $\ell$ por tipo de llamada, en el set de prueba con las marcas del proxy y el
modelo congelado. Corrida hecha para este informe.}\label{tab:ell}
\small
\begin{tabular}{@{}lrrrr@{}}
\toprule
Tipo de llamada & Llamadas & $\ell = -3$ & Intermedio & $\ell = +3$ \\
\midrule
\tabladatos{tab_ell}
\bottomrule
\end{tabular}
\end{table}

La tabla dice la regla. El modelo da $+3$ a las peticiones HTTP hacia un destino que no es destino
permitido de ninguna clave, sea o no un destino marcado como no confiable. A todo lo demás (leer
archivos, usar \cod{git}, hablar con la API de la empresa) le da $-3$, con pocas excepciones que
quedan en la columna «Intermedio» y que por sí solas no alcanzan ningún umbral.

Se entiende mirando las señales una por una. Una petición a un \esp{host} que no está en el grafo de
archivos tiene $m = 0$, $m_{\mathrm{sd}} = 0$ y \cod{jump} $= 0$, y esas tres cajas suman cerca de
$+3{,}3$. Que sea un sumidero suma $+4{,}6$. Con eso ya está sobre el tope cuando además desciende de
contenido no confiable ($+7{,}8$). Cuando no desciende ($-6{,}1$), el resultado lo decide la entropía
de los argumentos: en 269 de las 300 peticiones a \cod{attacker.test} del set llega igual a $+3$, y en
las otras 31 queda en \num{1.88}. Una petición a la API de la empresa, en cambio, toca un nodo del
grafo con sensibilidad alta, y las tres señales del grafo de archivos caen en cajas muy negativas: el
modelo aprendió que hablar con la API propia es trabajo normal.

\subsection{Qué significa esto}

\begin{itemize}
  \item \textbf{Por qué generaliza.} La regla no depende de la codificación ni del canal exacto: mira
    hacia dónde va la llamada. Por eso frena familias que usan base32 o XOR, que el contador no sabe
    leer.
  \item \textbf{Por qué rota de más.} Una sesión normal que lee una clave y después consulta una
    página web hace exactamente lo que la regla marca. Es la familia B5, donde Alicanto rotó de más
    en las 30 sesiones.
  \item \textbf{Por qué no ve la familia A9.} Los fragmentos viajan hacia la API de la empresa, que
    para el modelo es trabajo normal.
  \item \textbf{Cuánto aportan los grafos.} Según las ablaciones hechas después sobre validación
    (\cod{results/test/LECTURA.md}), quitar las señales de los dos grafos no baja los ataques frenados,
    pero sube las rotaciones en sesiones normales de 0,47 a 7,5 por cada \num{1000} llamadas. Un
    indicador binario simple (si la llamada toca algún archivo conocido) da lo mismo que la
    sensibilidad calculada. La diferencia de área bajo la curva al separar llamadas de ataque de
    normales, con y sin el grafo de archivos, es de \num{0.0045}, con intervalo al \qty{95}{\percent}
    de \num{0.0027} a \num{0.0065}: existe y es muy chica.
\end{itemize}
